% GENERATED FILE. Edit canonical lessons, then run npm run generate:books.
% canonical-source-hash: fnv1a64:23efd0d48bc147d2
% canonical-lessons: FR-C478-sourcil, FR-C478-paupiere, FR-C478-cil, FR-C478-narine, FR-C478-tempe

\chapter{Le visage de près}
\label{ch:fr-le-visage-de-pres}
\hlchaptermodality{478}

\begin{chapteropening}
\textbf{By the end of this chapter:} \emph{I can name the small parts of the face.}

Complete the last lesson of chapter 478: I can name the small parts of the face.
\end{chapteropening}

\section[le sourcil]{le sourcil — eyebrow}

\label{lesson:FR-C478-sourcil}



\begin{quote}
Before the new one: say the French for toe, then the French for hip.
\end{quote}

\subsection*{You'll want to know: le sourcil}
\textbf{le sourcil} — \textquotedblleft{}eyebrow\textquotedblright{}.

\begin{quote}
Il lève les sourcils quand il est surpris.
\end{quote}

\begin{quote}
\emph{He raises his eyebrows when he is surprised.}
\end{quote}

\emph{Froncer les sourcils} — to frown.

\begin{cousinweb}
\emph{le sourcil} is inherited from Latin \emph{supercilium} \textquotedblleft{}eyebrow\textquotedblright{} (EtymDB), built from \emph{super} \textquotedblleft{}above\textquotedblright{} and \emph{cilium} \textquotedblleft{}eyelid\textquotedblright{} (Lewis \& Short).
\end{cousinweb}

\begin{grammarlens}[title={What you've built}]
One more word for telling, explaining and describing at B1.
\end{grammarlens}

\subsection*{Your turn}
\begin{itemize}
\raggedright
  \item \emph{Say it:} \emph{le sourcil}
  \item \emph{Say it:} \emph{le sourcil}, once more
  \item \emph{Say it:} a sentence of your own with \emph{le sourcil}
  \item \emph{Recall:} say the French for toe, then the French for hip, then say \emph{le sourcil} again
\end{itemize}

\begin{tcolorbox}[breakable,colback=teal!4,colframe=teal!35!black,title={Before you move on}]
What does le sourcil mean? (\textquotedblleft{}Eyebrow\textquotedblright{}.) Say it once more.
\end{tcolorbox}
\section[la paupière]{la paupière — eyelid}

\label{lesson:FR-C478-paupiere}



\begin{quote}
Before the new one: say the French for hip, then the French for eyebrow.
\end{quote}

\subsection*{You'll want to know: la paupière}
\textbf{la paupière} — \textquotedblleft{}eyelid\textquotedblright{}.

\begin{quote}
Le soir, mes paupières deviennent lourdes.
\end{quote}

\begin{quote}
\emph{In the evening my eyelids get heavy.}
\end{quote}

\emph{Fermer les paupières} — to close one's eyes.

\begin{cousinweb}
\emph{la paupière} is inherited from Latin \emph{palpebra} \textquotedblleft{}eyelid\textquotedblright{} (Lewis \& Short), with a reshaped ending.
\end{cousinweb}

\begin{grammarlens}[title={What you've built}]
One more word for telling, explaining and describing at B1.
\end{grammarlens}

\subsection*{Your turn}
\begin{itemize}
\raggedright
  \item \emph{Say it:} \emph{la paupière}
  \item \emph{Say it:} \emph{la paupière}, once more
  \item \emph{Say it:} a sentence of your own with \emph{la paupière}
  \item \emph{Recall:} say the French for hip, then the French for eyebrow, then say \emph{la paupière} again
\end{itemize}

\begin{tcolorbox}[breakable,colback=teal!4,colframe=teal!35!black,title={Before you move on}]
What does la paupière mean? (\textquotedblleft{}Eyelid\textquotedblright{}.) Say it once more.
\end{tcolorbox}
\section[le cil]{le cil — eyelash}

\label{lesson:FR-C478-cil}



\begin{quote}
Before the new one: say the French for eyebrow, then the French for eyelid.
\end{quote}

\subsection*{You'll want to know: le cil}
\textbf{le cil} — \textquotedblleft{}eyelash\textquotedblright{}.

\begin{quote}
Le bébé a de longs cils noirs.
\end{quote}

\begin{quote}
\emph{The baby has long black eyelashes.}
\end{quote}

Mostly plural: \emph{les cils}.

\begin{cousinweb}
\emph{le cil} is inherited from Latin \emph{cilium} \textquotedblleft{}eyelid\textquotedblright{} (Lewis \& Short); in French it came to name the hairs on the eyelid.
\end{cousinweb}

\begin{grammarlens}[title={What you've built}]
One more word for telling, explaining and describing at B1.
\end{grammarlens}

\subsection*{Your turn}
\begin{itemize}
\raggedright
  \item \emph{Say it:} \emph{le cil}
  \item \emph{Say it:} \emph{le cil}, once more
  \item \emph{Say it:} a sentence of your own with \emph{le cil}
  \item \emph{Recall:} say the French for eyebrow, then the French for eyelid, then say \emph{le cil} again
\end{itemize}

\begin{tcolorbox}[breakable,colback=teal!4,colframe=teal!35!black,title={Before you move on}]
What does le cil mean? (\textquotedblleft{}Eyelash\textquotedblright{}.) Say it once more.
\end{tcolorbox}
\section[la narine]{la narine — nostril}

\label{lesson:FR-C478-narine}



\begin{quote}
Before the new one: say the French for eyelid, then the French for eyelash.
\end{quote}

\subsection*{You'll want to know: la narine}
\textbf{la narine} — \textquotedblleft{}nostril\textquotedblright{}.

\begin{quote}
Le cheval souffle par les narines.
\end{quote}

\begin{quote}
\emph{The horse is snorting through its nostrils.}
\end{quote}

Usually plural: \emph{les narines}.

\begin{cousinweb}
\emph{la narine} is inherited from Vulgar Latin \emph{nārīna} (EtymDB), built on Latin \emph{nāris} \textquotedblleft{}nostril\textquotedblright{} (Lewis \& Short).
\end{cousinweb}

\begin{grammarlens}[title={What you've built}]
One more word for telling, explaining and describing at B1.
\end{grammarlens}

\subsection*{Your turn}
\begin{itemize}
\raggedright
  \item \emph{Say it:} \emph{la narine}
  \item \emph{Say it:} \emph{la narine}, once more
  \item \emph{Say it:} a sentence of your own with \emph{la narine}
  \item \emph{Recall:} say the French for eyelid, then the French for eyelash, then say \emph{la narine} again
\end{itemize}

\begin{tcolorbox}[breakable,colback=teal!4,colframe=teal!35!black,title={Before you move on}]
What does la narine mean? (\textquotedblleft{}Nostril\textquotedblright{}.) Say it once more.
\end{tcolorbox}
\section[la tempe]{la tempe — temple (side of the head)}

\label{lesson:FR-C478-tempe}



\begin{quote}
Before the new one: say the French for eyelash, then the French for nostril.
\end{quote}

\subsection*{You'll want to know: la tempe}
\textbf{la tempe} — \textquotedblleft{}temple (side of the head)\textquotedblright{}.

\begin{quote}
Quand je suis fatigué, j'ai mal aux tempes.
\end{quote}

\begin{quote}
\emph{When I am tired, my temples ache.}
\end{quote}

A temple for worship is \emph{un temple}, a different word.

\begin{cousinweb}
\emph{la tempe} is inherited from Vulgar Latin \emph{tempula} (EtymDB), from Latin \emph{tempora} \textquotedblleft{}the temples of the head\textquotedblright{} (Lewis \& Short), a different word from \emph{tempus} \textquotedblleft{}time\textquotedblright{}.
\end{cousinweb}

\begin{grammarlens}[title={What you've built}]
One more word for telling, explaining and describing at B1.
\end{grammarlens}

\subsection*{Your turn}
\begin{itemize}
\raggedright
  \item \emph{Say it:} \emph{la tempe}
  \item \emph{Say it:} \emph{la tempe}, once more
  \item \emph{Say it:} a sentence of your own with \emph{la tempe}
  \item \emph{Recall:} say the French for eyelash, then the French for nostril, then say \emph{la tempe} again
\end{itemize}

\begin{tcolorbox}[breakable,colback=teal!4,colframe=teal!35!black,title={Before you move on}]
What does la tempe mean? (\textquotedblleft{}Temple (side of the head)\textquotedblright{}.) Say it once more.
\end{tcolorbox}
